% Source PDF: source/普通高考/2019/2019全国1文(河南,河北,山西,江西,湖北,湖南,广东,安徽,福建,山东).pdf, page 1.
% PDF embedded bitmap attempt: pdfimages -list found only an unrelated page image; this flowchart is vector content.
% Grid evidence: tmp/review_flowchart_2019_national_paper_1_liberal/grid/page-1-grid100.png
% (100px coarse + 50px fine) and q9-block-grid50.png (50px + 25px fine).
% Original comparison crop: tmp/review_flowchart_2019_national_paper_1_liberal/crop/q9_original_final.png,
% cropped from the no-grid page render.
% Elements: rounded start/end boxes, assignment rectangles, decision diamond, blank process rectangle,
% output parallelogram, directed connectors, loop-back branch and 是/否 labels.
% Relations: 开始 -> A=1/2 -> k=1 -> k<=2; 是 -> blank process -> k=k+1 -> loop back;
% 否 -> 输出 A -> 结束.  The decision and loop topology follows the source figure.
% solid segments: all box outlines and connector lines; dashed segments: none.
% labels: conditions centered in the diamond; 是 below the diamond, 否 above the right branch;
% math labels centered in boxes, with the loop-back path kept clear of text.
% Directed edges: 开始→A=1/2→k=1→判定；判定(是)→空白框→k=k+1→回到判定入口；
% 判定(否)→输出A→结束。回路独立沿左侧绕行，不与输出分支共享线段或假交汇。
\begin{tikzpicture}[
  >=Stealth,
  line width=0.45pt,
  line join=round,
  line cap=round,
  font=\normalsize,
  flow/.style={draw, line width=0.45pt, anchor=center, align=center, inner sep=2pt,
               execute at begin node={\setlength{\parindent}{0pt}}},
  startstop/.style={flow, rounded corners=0.18cm, minimum width=1.05cm, minimum height=0.52cm},
  process/.style={flow, rectangle, minimum width=1.80cm, minimum height=0.82cm},
  smallprocess/.style={flow, rectangle, minimum width=1.15cm, minimum height=0.52cm},
  decision/.style={flow, diamond, aspect=2.9, minimum width=3.05cm, minimum height=1.10cm, inner xsep=0pt},
  output/.style={flow, trapezium, trapezium left angle=72, trapezium right angle=108,
                minimum width=1.55cm, minimum height=0.72cm},
]
  % Schematic coordinates preserve the source topology and relative branch placement.
  \node[startstop] (start) at (0,6.85) {开始};
  \node[process]   (init)  at (0,5.62) {$A=\dfrac{1}{2}$};
  \node[smallprocess] (kinit) at (0,4.45) {$k=1$};
  \node[decision] (test) at (0,3.22) {$k\leqslant 2$};
  \node[process] (blank) at (0,1.88) {};
  \node[process] (inc) at (0,0.62) {$k=k+1$};
  \node[output] (out) at (2.65,1.88) {输出 $A$};
  \node[startstop] (finish) at (2.65,0.62) {结束};

  \draw[->] (start.south) -- (init.north);
  \draw[->] (init.south) -- (kinit.north);
  \draw[->] (kinit.south) -- (test.north);
  \draw[->] (test.south) -- node[right=2pt] {是} (blank.north);
  \draw[->] (blank.south) -- (inc.north);
  \draw[->] (test.east) -- node[above=2pt] {否} (2.65,3.22) -- (2.65,2.24);
  \draw[->] (out.south) -- (finish.north);

  % Loop-back from the increment box to the decision's incoming node.
  \draw[->] (inc.west) -- (-1.75,0.62) -- (-1.75,3.75) -- (-0.35,3.75) -- (0,3.75);
\end{tikzpicture}
